% ============================================================================
% METHODOLOGY UPDATE -- brings Section III in line with the two experiments
% actually reported in Section IV.
%
% Four edits, marked EDIT 1..4. Each gives the text to replace and the text
% to replace it with. Nothing outside the Methodology section changes.
%
% WHY: the Methodology currently documents one protocol (the 2,000-run sweep,
% seeds 10/20/30/40/50, cudnn.deterministic=True). Section IV-F reports a
% second experiment run under a different seed set and a different seeding
% scope, which the Methodology never describes and, in two places, actively
% contradicts.
% ============================================================================


% ---------------------------------------------------------------------------
% EDIT 1 -- Training Configuration and Evaluation
% Scope the "all runs" claim to the primary sweep.
%
% REPLACE:
%   Every setting in Table~\ref{tab:hyperparams} was held the
%   same across all runs, so the activation pair is the only
%   thing that changes between models.
%
% WITH:
% ---------------------------------------------------------------------------

Every setting in Table~\ref{tab:hyperparams} was held the same
across the 2{,}000 runs of the primary sweep, so the activation
pair is the only thing that changes between models. The
early-learning follow-up of Section~\ref{sec:convergence_protocol}
departs from this in three documented respects.


% ---------------------------------------------------------------------------
% EDIT 2 -- Replication Protocol
% Name the sweep, and record cudnn.benchmark, which the follow-up differs on.
%
% REPLACE the first paragraph of \subsection{Replication Protocol}
% ("Each of the $10\times10\times4 = 400$ configurations ... exactly
% reproducible.") WITH:
% ---------------------------------------------------------------------------

In the primary sweep, each of the $10\times10\times4 = 400$
configurations was trained under five seeds,
$\{10,20,30,40,50\}$ (Table~\ref{tab:hyperparams}), giving
$2{,}000$ models. The seed is applied immediately before the
model and its data loaders are constructed, setting
\texttt{random}, \texttt{numpy}, \texttt{torch} and
\texttt{torch.cuda}, with \texttt{cudnn.deterministic=True} and
\texttt{cudnn.benchmark=False} to pin kernel selection as well.
A given (pair, seed) is therefore exactly reproducible.


% ---------------------------------------------------------------------------
% EDIT 3 -- NEW subsection.
% INSERT immediately after \subsection{Replication Protocol} and before
% \subsection{On Reproducing the Baseline of Ahmed et al.}
% ---------------------------------------------------------------------------

\subsection{Early-Learning Follow-Up Protocol}
\label{sec:convergence_protocol}

The early-learning experiment reported in
Section~\ref{sec:learning} is a separate, smaller study. It does
not inherit the primary sweep's protocol, and we state the
differences here because two of them affect how its numbers
should be read.

\textit{Scope.} It trains BiGRU+GRU only, and only the ten
configurations that use the same function in both dense layers
(X+X), rather than all 100 ordered pairs. The dataset, split,
vocabulary, architecture sizes, optimiser, scheduler, batch size,
maximum epochs and early-stopping rule are all as in
Table~\ref{tab:hyperparams}.

\textit{Measurement.} It records validation loss and validation
accuracy after every epoch. The test set is never touched. Its
accuracies are therefore validation accuracies on the held-out
10\% of the training split, and are not comparable with the test
accuracies reported elsewhere in this paper.

\textit{Seeding.} This experiment sets \texttt{torch.manual\_seed}
and \texttt{numpy.random.seed} only, and leaves the cuDNN flags at
their library defaults. The narrower scope is deliberate: it keeps
the runs bit-compatible with the earlier single-seed convergence
result, and enabling \texttt{cudnn.deterministic} changes the
numerics enough to flip epoch-one outcomes. That sensitivity is
itself a finding. Epoch one is close to a threshold, in which a
model either leaves the chance-level plateau within one pass or
does not, so small numerical differences can decide the epoch-one
figure while leaving later epochs unaffected. These runs are
reproducible on a fixed machine and library version, but we do not
claim reproducibility across different cuDNN kernel selections.

\textit{Seed sets.} The follow-up was run twice, under two
disjoint five-seed sets: $\{42,0,1,2,3\}$ and
$\{10,20,30,40,50\}$, for 100 runs in total. The first set
contains seed 42, which is also the seed of the single-run
ranking it is compared against, so on its own it cannot serve as
an independent check of that ranking. The second set shares no
seed with the first and matches the seeds of the primary sweep,
which allows the epoch-one ranking to be tested on runs that have
no seed in common with those used to establish it.


% ---------------------------------------------------------------------------
% EDIT 4 -- Table~\ref{tab:hyperparams}
%
% REPLACE the row:
%   Training seeds                & 10, 20, 30, 40, 50 & Ours \\
%
% WITH the two rows:
% ---------------------------------------------------------------------------

    Training seeds (primary sweep) & 10, 20, 30, 40, 50 & Ours \\
    Training seeds (follow-up)    & $\{42,0,1,2,3\}$, $\{10,\ldots,50\}$ & Ours \\
